% ITI 2026 -- NIT Agartala, 13-14 November 2026. Springer LNNS, 10-12 pages.
%
% Build:  pdflatex iti2026 && bibtex iti2026 && pdflatex iti2026 && pdflatex iti2026
%
% Written to the ITI 2026 rules:
%   * required sections: abstract, keywords, introduction, literature review,
%     methodology, results, discussion, conclusion, references
%   * numbered citations in brackets, author surnames named in the text
%   * every table typed (no images), every equation typed (no images)
%   * every table referred to and discussed in the text
%   * LLM/tool use disclosed in the methodology (Section 3.5), as the rules require
%
% Style: short sentences, plain words, active voice. One idea per sentence.

\documentclass[runningheads]{llncs}

\usepackage[T1]{fontenc}
\usepackage{booktabs}
\usepackage{amsmath}
\usepackage{amssymb}

\begin{document}

\title{Pricing Deception: Deciding When a Web Attack\\Detector Should Probe}
\titlerunning{Pricing Deception}

\author{Parth Kansal \and Amrit Singh Nijjar \and Siddhant Mehta}
\authorrunning{P. Kansal et al.}
\institute{Department of AIT-CSE, Chandigarh University,\\
Gharuan, Mohali, Punjab, India\\
\email{parth.kansal823@gmail.com, amrit.nijj814@gmail.com, 23bis70162@cuchd.in}}

\maketitle

% ---------------------------------------------------------------- abstract ~120w
\begin{abstract}
A web application firewall must decide to block or allow each request using only the
information that request contains. Deception is normally used later, after a session
has already been judged hostile. We use it earlier, while the decision is still open.
When the system is unsure about a client, it adds a small invisible probe to the
response. A browser never shows the probe on the page, while a client that acts on it
reveals itself. The system decides when to probe by computing what the answer is
worth, not by a threshold set by hand. Under ordinary cost accounting no middle
option exists at all, so the probe cannot be tuned into existence. Over 99 runs,
probing caused an extra 7.0\% of attackers to be caught, and no legitimate session
was ever diverted.

\keywords{Cyber deception \and Intrusion detection \and Web application security
\and Value of information \and Cost-sensitive learning \and Honeytokens}
\end{abstract}

% ------------------------------------------------------------ 1 introduction
\section{Introduction}

Building a probe is engineering. The question we answer is \emph{when} the detector
should use one, and our answer is that this can be calculated. A probe gives no
immediate protection, because the request still reaches the real application; its
only value is the information it may reveal. Priced by that value
\cite{howard1966information}, probing becomes the cheapest action over a band of
belief whose edges follow from a fixed cost table and a measured probe success rate,
not from our judgement.

\paragraph{Contributions.} Two.

\begin{enumerate}
\item \textbf{A priced third action.} Probing is the cheapest action over a band of
belief whose edges are calculated, not tuned. Under ordinary cost accounting that
band is empty, so there is no middle option to tune: the probe appears only once the
value of information is included.

\item \textbf{A measured causal effect.} Rather than compare a system with probes
against one without, we withhold the probe from a random one in ten sessions inside
the probing system. Both groups sit at the same belief, so the gap between them
measures what the probe itself caused.
\end{enumerate}

% -------------------------------------------------------- 2 literature review
\section{Literature Review}

We group earlier work into four areas and state what each leaves open.

\paragraph{Detecting web attacks.} Signature firewalls match known bad patterns, so
they are easy to defeat by obfuscation and blind to attacks using valid syntax.
Learned detectors replace fixed patterns with a model: Kruegel and Vigna
\cite{kruegel2003anomaly} built an early anomaly detector for web requests,
Robertson et al. \cite{robertson2006generalization} generalised it across sites, and
later work applies deep \cite{tekerek2021novel} and adaptive \cite{amouei2022rat}
methods. All share one habit. They watch a session gather evidence and then commit
to a score. Our detector belongs to this family and is not our contribution; we use
a signature firewall and a passive detector as baselines, because our contribution
sits on top of them.

\paragraph{Telling scripts from people.} A separate line asks whether a client is
automated, not whether it is hostile. Iliou et al. \cite{iliou2021detection} show
that advanced bots reproduce browser behaviour closely enough that automation
signals alone stop separating them from people. This is why our detector scores
automation and hostility separately. A scanner is automated and hostile; a
price-comparison bot is automated and harmless; a careful human attacker is neither.
One score cannot express the middle case.

\paragraph{Deception and honeytokens.} Deception as a defence is well studied.
Provos \cite{provos2004honeyd} built an early virtual-honeypot framework, Almeshekah
and Spafford \cite{almeshekah2014planning} argue deception must be planned around a
target belief, and Pawlick et al. \cite{pawlick2019taxonomy} give a taxonomy.
Beltr\'an-L\'opez et al. \cite{beltran2025cyberdeception} survey the current state of
the art, and Ferguson-Walter et al. \cite{fergusonwalter2021examining} study whether
decoys change what attackers do. In nearly all of it the deception is a \emph{place}:
a caught attacker is moved there, after the decision has been made. Planting a fake
file or password as a tripwire is equally long established -- honeyfiles
\cite{yuill2004honeyfiles}, decoy documents \cite{bowen2009baiting} and fake
passwords \cite{juels2013honeywords} -- but that work treats traps as always on.
Kahlhofer and Rass \cite{kahlhofer2024applayer} review deception placed inside the
application itself and Kahlhofer et al. \cite{kahlhofer2025koney} automate deploying
and rotating these traps, yet neither decides \emph{when} to place one from what the
system currently believes. Traps are also not free to hide, having been fingerprinted
at scale for both honeypots \cite{vetterl2018bitter} and honeytokens
\cite{srinivasa2020honeytoken}.

\paragraph{Generated fake environments.} Language models increasingly produce
convincing fake systems. Sladi\'c et al. \cite{sladic2024shellm} generate shell
output live, Bridges et al. \cite{bridges2025sok} review the area and find the gains
so far modest, and honeypots are now themselves monitored for language-model
attackers \cite{reworr2024llmagent}. What this work measures is realism: how
tempting a planted item looks \cite{kahlhofer2024honeyquest}, and how realistic a
generated one reads \cite{timmer2025honeyfile}. Vero et al. \cite{vero2026honeyval}
propose an evaluation framework covering stealth and fidelity, but not whether a
fake system contradicts itself during a long visit. Our fake site is built so its
content generator can be swapped, and we test that separately.

\paragraph{Decision theory.} The mathematics we use is standard. Howard
\cite{howard1966information} defined the value of sample information, Elkan
\cite{elkan2001foundations} set out cost-sensitive decisions under a loss matrix,
Pawlick et al. \cite{pawlick2019leaky} model deception as a signalling game with
evidence, and Axelsson \cite{axelsson2000baserate} showed how the rarity of attacks
constrains any intrusion detector. None of this includes an action that buys
evidence. That is the gap we fill. We claim no new mathematics: our contribution is
noticing that a probe \emph{is} a purchase of information, and that pricing it as
one turns a hand-set option into a calculated one. We follow the evaluation warnings
of Sommer and Paxson \cite{sommer2010outside} and Arp et al. \cite{arp2022dos}
throughout.

% ------------------------------------------------------------ 3 methodology
\section{Methodology}

\subsection{System Overview}

Our system is a reverse proxy. It sits in front of the real web application.
Clients only ever talk to the proxy. The application never sees a request that has
not passed through it, and never knows the deception exists.

Every request follows the same five steps.

\begin{enumerate}
\item Identify the session, using a cookie.
\item Turn the request into 18 numbers, called features.
\item Update two running scores for the session.
\item Price the three possible actions and take the cheapest.
\item Write the decision to a log that cannot be edited afterwards.
\end{enumerate}

The three actions are \emph{pass}, \emph{probe} and \emph{divert}. Pass sends the
request to the real application. Divert moves the session into a fake copy of the
site. Probe adds an invisible probe to the response and passes the request on
unchanged.

\subsection{Two Scores, Not One}

The 18 features split into two groups, because the threat has two parts.

Ten \emph{automation} features describe how the client behaves. They include the
time between requests, how regular that timing is, whether images and stylesheets
were fetched, how complete the browser headers are, and whether a cookie was sent.

Eight \emph{hostility} features describe what the client is trying to do. They
include how many unusual characters appear in user input, whether database
keywords appear in a suspicious position, how many replies were errors, how many
logins failed, and how many different accounts the session tried to log in to.

Two logistic models turn these features into two scores. Both scores build up
across the session. We keep them apart for a reason. The cost of an action depends
only on hostility, because an automated client may be harmless. So hostility drives
the decision. Automation is used only to choose \emph{which} probe to use, because
a script and a human take different bait.

\subsection{Choosing an Action by Price}

At any moment the system holds a belief $p$ that the visitor is hostile. Each action
has a cost under a loss matrix in the sense of Elkan \cite{elkan2001foundations}, and
those costs are fixed and locked before any data is collected.
Table~\ref{tab:costs} gives them.

\begin{table}[t]
\centering
\caption{The fixed cost table. These numbers were locked and hashed before any
data was collected, so no threshold below could be adjusted to improve results.}
\label{tab:costs}
\begin{tabular}{@{}lrrr@{}}
\toprule
 & pass & probe & divert \\
\midrule
visitor is harmless & 0  & 1  & 200 \\
visitor is hostile  & 25 & 25 & $-20$ \\
\bottomrule
\end{tabular}
\end{table}

Three numbers in Table~\ref{tab:costs} carry the argument. Diverting a harmless
visitor is the worst outcome, so it costs 200. Probing a hostile visitor costs the
same as passing one, both 25, because the request still reaches the real
application either way. The only extra cost of probing is 1, paid by harmless
visitors.

Using these costs, the expected cost of each action is a straight line in $p$:
%
\begin{equation}
\mathbb{E}[C_{\text{pass}}] = 25p, \qquad
\mathbb{E}[C_{\text{probe}}] = 1 + 24p, \qquad
\mathbb{E}[C_{\text{divert}}] = 200 - 220p .
\label{eq:costs}
\end{equation}

Now compare the first two lines of Equation~(\ref{eq:costs}):
%
\begin{equation}
\mathbb{E}[C_{\text{probe}}] - \mathbb{E}[C_{\text{pass}}] = 1 - p \;\geq\; 0 .
\label{eq:gap}
\end{equation}

Equation~(\ref{eq:gap}) is never negative. Probing always costs more than passing.
So on cost alone the system would never probe. It would only pass or divert, and
switch between them where $25p = 200 - 220p$, that is at $p \approx 0.816$.

This matters more than it may appear. It means there is no middle option to tune.
A reader may reasonably suspect that our probing band is just a threshold we chose.
It is not, because without the next step there is no band at all.

\subsection{Pricing the Information}

A probe is worth something only because of what it may tell us. It produces one of
two outcomes: the client takes the bait, or it does not. Either way the belief
becomes sharper, and the next decision becomes better.

We measure this using the expected value of sample information \cite{howard1966information}:
%
\begin{equation}
V(p) = \min_a \mathbb{E}[C(a) \mid p] \;-\;
       \mathbb{E}_Z\!\left[ \min_a \mathbb{E}[C(a) \mid p'(Z)] \right].
\label{eq:evsi}
\end{equation}

In Equation~(\ref{eq:evsi}), $Z$ is what we observe, and $p'(Z)$ is the updated
belief after seeing it. Let $\beta_a$ be the chance a hostile client takes the
bait, and $\beta_b$ the chance a harmless one does. Then the chance of a bite is
$p\beta_a + (1-p)\beta_b$, and the updated beliefs follow from Bayes' rule. Every
term inside Equation~(\ref{eq:evsi}) is a straight line, so $V(p)$ can be worked
out directly. No simulation is needed.

The system then prices probing at its own cost minus the value of the information
it buys, and picks the cheapest action:
%
\begin{equation}
\text{price}(\text{probe}) = 1 + 24p - V(p).
\label{eq:price}
\end{equation}

Two simple facts fix the band at both ends. First, $V(p) \ge 0$ always, because
the smallest of several straight lines is a concave function, and information can
never make a decision worse. Second, $V(0) = V(1) = 0$, because when we are already
certain, no answer can change what we do. So the band cannot spread across the
whole range, and it cannot be widened by tuning.

With our fixed costs and measured bite rates, the calculated result is:
%
\begin{equation}
\text{pass if } p < 0.065, \qquad
\text{probe if } 0.065 \le p < 0.879, \qquad
\text{divert if } p \ge 0.879 .
\label{eq:bands}
\end{equation}

One detail in Equation~(\ref{eq:bands}) is worth noticing. Offering a probe pushes
the divert point \emph{up}, from 0.816 to 0.879. A system that can probe is more
patient with an honest user than one that cannot. We did not aim for this. It
falls out of the arithmetic.

\subsection{Measuring the Probe, and Tools Used}
\label{sec:tools}

The calculation above needs one measured input: how often each kind of visitor
takes the bait. We measure this in a separate round, run after the probes were
built and locked before the evaluation. For our main probe, 138 of 183 hostile
sessions took the bait, and 0 of 74 harmless sessions did. This gives
$\beta_a = 0.753$ and $\beta_b = 0.0067$. We do not set $\beta_b$ to zero, even
though no harmless session took the bait, because that would make the arithmetic
infinite and would assume away the safety property we are trying to measure.

As the conference guidelines require, we state which tools were used and where. All
detection code, features, decision rule and experiments are our own, written in
Python; no language model is used anywhere in the request path, nor for data
processing or filtering, visualisation, or proving any result. Language models were
used in exactly three places, all of which affect results reported here. First, as
\textbf{simulated attackers}: three local open models (Llama 3.2 1B, Llama 3.2 3B and
Qwen 2.5 7B) each read the replies and chose their own next request, none told that
probes existed, which tests our assumption about how curious an attacker is. Second,
as a \textbf{content generator} for the fake site in one comparison, to check that
our consistency layer does not depend on how the content is produced. Third, for
\textbf{adversarial testing}, attacking our own running system to look for mistakes;
it found several flaws in the fake site, which we fixed and report in
Section~\ref{sec:discussion}.

\subsection{Experimental Setup}

We compare three systems. They differ only in one setting, and all use the same
locked model.

\begin{itemize}
\item \textbf{B1}, a signature firewall using fixed patterns.
\item \textbf{B2}, our detector without probes. It can only pass or divert.
\item \textbf{B4}, the full system, with probes and the fake site.
\end{itemize}

Each system was run over 99 separate rounds of generated traffic. Each round has
120 attack sessions and 80 harmless ones. That gives 11{,}880 attack sessions and
7{,}920 harmless sessions per system. Within one round, all three systems see
exactly the same traffic, so every attack session can be compared directly across
systems. We fixed the main comparison before running anything.

Throughout this paper, \emph{diverted} means the system decided the visitor was
hostile. It is counted the same way in every system, so the numbers are comparable.
The design follows the evaluation guidance for machine learning in security
\cite{sommer2010outside,arp2022dos}: the comparison is paired and pre-registered, and
the attack rounds that train and test the detector are disjoint.

% --------------------------------------------------------------- 4 results
\section{Results}
\label{sec:results}

\subsection{What the Probe Itself Causes}

Comparing B2 and B4 compares two whole systems, which differ in more than the
probe. To measure the probe alone, we withhold it from about one session in ten
among those that reach the probing band. The choice is random. Both groups sit at
the same belief, under the same rules. The only difference is whether a probe was
sent. Table~\ref{tab:holdout} gives the outcome.

\begin{table}[t]
\centering
\caption{The randomised comparison inside the full system. Both groups reached the
probing band under identical rules. Only the probe differs.}
\label{tab:holdout}
\begin{tabular}{@{}lrrr@{}}
\toprule
group & sessions & caught & rate \\
\midrule
probe sent      & 10{,}643 & 10{,}112 & 0.950 \\
probe withheld  & 1{,}237  & 1{,}089  & 0.880 \\
\bottomrule
\end{tabular}
\end{table}

Table~\ref{tab:holdout} shows a difference of \textbf{7.0 percentage points}. The
95\% confidence interval is $[5.2, 8.8]$, and Fisher's exact test gives
$p = 3.4 \times 10^{-19}$. Because the two groups were separated at random at the
same belief, this difference is caused by the probe. It is not a comparison of two
different designs.

We also checked the random split was fair: the largest compositional difference
between the groups across the five attack types is 1.2 percentage points, and a
chi-square test gives 2.58 on four degrees of freedom, well below the 9.49 that
would signal a problem. An unfair split would have measured the mixture, not the
probe.

\subsection{Comparison With Baselines}

Table~\ref{tab:baselines} compares the three systems.

\begin{table}[t]
\centering
\caption{Attack recall and harm to normal users, over 99 rounds. Higher recall is
better. For diverted harmless sessions, zero is best.}
\label{tab:baselines}
\begin{tabular}{@{}llrr@{}}
\toprule
system & attack recall (95\% CI) & spread & harmless diverted \\
\midrule
B1 signature firewall & 0.366 [0.358, 0.375] & 0.020 & 0 / 7{,}920 \\
B2 detector only      & 0.889 [0.883, 0.894] & 0.024 & 4 / 7{,}920 \\
B4 full system        & \textbf{0.943 [0.939, 0.947]} & 0.018 & \textbf{0 / 7{,}920} \\
\bottomrule
\end{tabular}
\end{table}

The two columns point in opposite directions -- higher recall is better, zero
diverted harmless sessions is best -- and B4 wins on both. The recall ranges for B2
and B4 do not overlap. Comparing the same sessions in both systems, 842 attacks were
caught by B4 but missed by B2, while 197 went the other way; a paired test gives
$p = 1.9 \times 10^{-95}$, and B4 led in all 99 rounds, which rules out a lucky
result.

Those 197 are not noise, and we do not present them as such. They are delays. Their
belief landed between 0.816 and 0.879, where B2 diverts at once and B4 first asks a
question, and some sessions end before the question is answered. That is the real
cost of having a third action.

The two zeros in the last column deserve an explanation. B1's is unsurprising: its
patterns need suspicious syntax that normal traffic does not contain. B4's is more
interesting, because B4 is not merely as safe as B2 but safer. All four harmless
sessions B2 wrongly diverted reached the probing band in B4 and were probed instead.
None took the bait, being genuinely harmless, and that silence kept their belief
below 0.879 where B2 had already acted at 0.816.

\subsection{Comparison With a Real Firewall}

B1 is a firewall we wrote, and a reader may reasonably distrust it, so we replayed
the same traffic through the OWASP ModSecurity Core Rule Set that many real sites
use. Table~\ref{tab:crs} gives the result at each of its four settings.

\begin{table}[t]
\centering
\caption{OWASP ModSecurity Core Rule Set on the same traffic. Perfect recall is
only reached at a setting no real site could use.}
\label{tab:crs}
\begin{tabular}{@{}lll@{}}
\toprule
setting & attack recall (95\% CI) & harmless sessions blocked \\
\midrule
1 (default) & 0.353 [0.305, 0.404] & 0 / 240 \\
2           & 0.544 [0.493, 0.595] & 0 / 240 \\
3           & 0.544 [0.493, 0.595] & 0 / 240 \\
4 (maximum) & \textbf{1.000} [0.989, 1.000] & \textbf{72 / 240 (30\%)} \\
\bottomrule
\end{tabular}
\end{table}

Two points follow. At its default setting the real rule set scores 0.353, close to
our B1 at 0.366, so B1 is not a weakened stand-in. And the limit is not a matter of
tuning: the rule set reaches perfect recall only by blocking 30\% of legitimate
visitors, which no real site could accept, and at every setting that leaves normal
users alone it stops near 0.54. Our full system reaches 0.943 while diverting
nobody. One point favours the rule set, and we state it plainly: it is judged on
blocking any single request, while our systems are judged on identifying the whole
session, which is the harder task.

\subsection{Where the Probe Helps}

Our theory predicts something specific. The probe should help where the detector is
genuinely unsure, and nowhere else. Table~\ref{tab:where} tests this.

\begin{table}[t]
\centering
\caption{Recall by attack type. The gain is concentrated in one type: the case
where the detector cannot decide on its own.}
\label{tab:where}
\begin{tabular}{@{}lrrrr@{}}
\toprule
attack type & sessions & B2 & B4 & gain \\
\midrule
Obfuscated SQL injection & 3{,}960 & 0.898 & 0.921 & $+0.022$ \\
Slow SQL injection       & 1{,}980 & 1.000 & 1.000 & $0.000$ \\
Record walking via API   & 1{,}980 & 0.994 & 0.993 & $-0.001$ \\
Password spraying        & 1{,}980 & 1.000 & 1.000 & $0.000$ \\
\textbf{Record walking via web pages} & 1{,}980 & \textbf{0.540} & \textbf{0.823} & $\mathbf{+0.282}$ \\
\bottomrule
\end{tabular}
\end{table}

Table~\ref{tab:where} shows the whole gain sitting in one row. Walking through
records using the normal web pages improves from 0.540 to 0.823. Everywhere else
the probe changes almost nothing. This is exactly the case a signature cannot see
and the detector cannot settle, which is where our theory says information is worth
buying.

One attack type cannot be helped at all, and it is half of what the full system
still misses. Obfuscated SQL injection never takes the bait, though nine in ten of
those sessions were shown one: those attackers send payloads without reading the
replies, so a probe placed in a reply cannot reach them.

\subsection{Do People Notice the Fake Site?}
\label{sec:human}

Diverting a visitor is only safe if the fake site does not announce itself, and
until now that judgement was our own. Each participant in a small blind study was
assigned at random either to the genuine application or to the fake copy, explored
it for about ten minutes with an ordinary staff login, and was then asked whether
the system was real or a mock-up. Both arms got the same instructions and the same
questions, nobody was told what the study was about, and assignment is
block-randomised from a fixed seed so it cannot be quietly re-drawn.

Four people took part, two per arm. One of the two shown the fake site called it a
mock-up; so did one of the two shown the genuine application. The difference between
the arms is zero, and Fisher's exact test gives $p = 1.0$.

That $p$ proves nothing on its own, and we would rather say so than let it flatter
us. With two participants per arm the test cannot return anything below $0.333$,
even had every participant in one arm and none in the other called the system fake.
No significant result was reachable at this size; our protocol asks for eight, the
smallest number at which a clean separation would reach $p < 0.05$. Showing that the
two arms are alike, rather than merely failing to separate them, is an equivalence
question \cite{schuirmann1987comparison} and needs a sample we do not yet have.

The study still supplies something no automated check can. A participant exploring
the \emph{genuine} application called it a mock-up after fifteen minutes, because
the staff directory grouped three departments together and left a fourth apart.
People call real systems fake on the strength of an ordinary layout quirk, so the
baseline this question must be read against is not zero.

% ------------------------------------------------------------ 5 discussion
\section{Discussion}
\label{sec:discussion}

Our results support the mechanism, not just a headline number. Had the probe raised
recall everywhere, that would suggest a weak detector rather than a correct
calculation. Instead the gain sits almost entirely in the one case the detector
cannot settle, which is what the pricing predicts: information is worth buying
exactly where the decision is in doubt, and worth nothing elsewhere. The safety
result has the same shape. Raising the divert point from 0.816 to 0.879 was not a
feature we added but a consequence of the arithmetic, and its effect was to recover
every mistake the simpler detector made.

\paragraph{Where our approach does not win.} Setting the band edges by hand beats
calculating them, if expected cost is the only measure, and we report that rather
than hide it. The belief takes only about two distinct values in our system, so the
extra precision of a calculated edge does no real work, and every setting that beats
it does so by diverting harmless users. Among settings that divert nobody, our
calculated edges are the best available: a smaller claim than we first hoped to
make, and a more useful one.

\paragraph{Adversarial testing of the fake site.} A language model agent attacking
our own system (Section~\ref{sec:tools}) found four ways to tell the fake site from
the real one, each found by attacking the fix for the one before: a page returning
different values after a move, a list page naming a person differently from their
own profile, and a record existing after a move that did not exist before. A human
volunteer later found a fifth, two profiles carrying the same name but different
details. All five are fixed. The lesson generalises: checking that a fake site agrees
with itself is not the same as checking that it agrees with what the attacker already
saw. This is a different failure from the fingerprinting of deployed decoys
\cite{vetterl2018bitter,srinivasa2020honeytoken}, which reads the artefact rather
than its history.

\paragraph{Who actually takes the bait.} Bite rate does not rise steadily with the
attacker's capability, and that bounds our central claim. Across the three local
models it climbs from 0.077 to 0.556 as the model gets stronger, then falls to zero
at both ends of the range. Off-the-shelf tools never read the replies. At the other
end, neither the frontier-model agent nor any of our hand-driven testers has taken a
single bait: four interactive sessions, thirty-four probes shown, none touched. All
four were still caught, by the injection path rather than by a probe. Our 7.0\% is
therefore corroborated only in the middle of the capability range. A careful reader
treats a planted database error as we would, as something to leave alone. How
enticing a planted item looks is itself measurable \cite{kahlhofer2024honeyquest},
and calibrating our probes against a human panel is the obvious next step. This
narrows where the mechanism applies rather than withdrawing it.

\paragraph{Limitations.} Four, in order of importance. First, all our traffic is
generated: every rate describes our own traffic and would need measuring again on
real traffic, and while the calculation and the randomised design do not depend on
the traffic being real, the numbers do. Second, we tune and test on one application;
the locked detector kept working in front of a second, very different one, but the
probes themselves would need rewriting. Third, four people have tested the fake site
(Section~\ref{sec:human}), too few for any arm difference to reach significance.
Fourth, the bite rate depends on how curious we assumed attackers are, and no human
we have observed has yet taken a probe.

% ------------------------------------------------------------ 6 conclusion
\section{Conclusion}

Passive defences force a choice between acting too early and acting too late, and
that choice can be avoided. A defender can create evidence by placing a small
invisible probe in a reply, and when to do so can be calculated rather than guessed.
Priced by the value of the information it buys, probing is the cheapest action over
a band that does not exist under ordinary cost accounting, so the probe cannot be a
tuned threshold. Our measurements support the mechanism: the gain appears only where
the detector is genuinely unsure, a randomised comparison shows the probe caused it,
and no legitimate session was ever diverted. What remains is to test the idea
against real traffic and real attackers.

\bibliographystyle{splncs04}
\bibliography{refs}

\end{document}
